% DME_functor.tex --- Diamond Master Equation as a partial functor
% Standalone LaTeX statement, math.OA / math.CT level.
% NO new theorem; see DME_functor.md for verdict and limits-of-scope.

\documentclass[11pt]{article}
\usepackage{amsmath,amssymb,amsthm}
\usepackage[margin=1in]{geometry}
\usepackage{tikz-cd}

\newtheorem{theorem}{Theorem}
\newtheorem{proposition}[theorem]{Proposition}
\newtheorem{definition}[theorem]{Definition}
\newtheorem{remark}[theorem]{Remark}
\newtheorem{question}[theorem]{Open Question}

\newcommand{\AA}{\mathcal{A}}
\newcommand{\NN}{\mathcal{N}}
\newcommand{\FRW}{{\mathrm{FRW}}}
\newcommand{\HL}{{\mathrm{HL}}}
\newcommand{\R}{\mathbb{R}}
\newcommand{\Rprop}{R_{\mathrm{proper}}}
\newcommand{\Diam}{\mathbf{Diam}_\FRW}
\newcommand{\WCatinj}{\mathbf{W^{*}\!Cat}_{\mathrm{inj}}}

\title{The Diamond Master Equation as a partial functor on the poset of admissible FRW comoving diamonds}
\author{(internal categorical companion to \texttt{frw\_note} + \texttt{krylov\_diameter})}
\date{2026-05-02}

\begin{document}
\maketitle

\begin{abstract}
We formalise the assignment $D\mapsto(\NN(D),K(D),C(D;s))$ of
\texttt{unifying\_equation\_attempt.md} (Candidate~C, the
\emph{Diamond Master Equation}) as a covariant assignment from the
poset $\Diam$ of admissible doubly-bounded comoving causal diamonds
in radiation-dominated FRW with $\eta_i>0$. We show that
$D\mapsto\NN(D)$ is a covariant functor into a Brunetti--Fredenhagen--Verch--style
target category $\WCatinj$ of W$^*$-algebras with normal injective
unital $*$-homomorphisms, conditional on the modular-compatibility
hypothesis of Takesaki's conditional-expectation theorem. The
modular Hamiltonian $K(D)$ does \emph{not} extend to an
operator-valued natural transformation (only its automorphism group
$\sigma^D$ does, by Connes 1973 cocycle invariance). The
Krylov--Diameter Correspondence (\texttt{krylov\_diameter} Thm.~4)
is a candidate natural transformation
$\phi:F_C\Rightarrow F_R^{-1}$ between two functors on $\Diam$;
we show that the naturality square fails on inclusions in general,
because the Lanczos coefficients of the spread complexity are not
preserved by the Connes--Takesaki conditional expectation
$E:\NN(D)\to\NN(D')$. The conclusion is a \emph{partial functor with
a non-natural transformation}, not a categorical equation in the
sense of Doplicher--Roberts or Brunetti--Fredenhagen--Verch. No new
theorem is proved.
\end{abstract}

\section{The source category $\Diam$}\label{sec:source}

\begin{definition}[Admissible diamond]\label{def:admissible}
Fix the radiation-dominated FRW background $a(\eta)=a_0\eta$ on
$\eta\in(0,\infty)$ (resp.\ matter $a=a_0\eta^2$ or de~Sitter
$a=-1/(H\eta)$ on $\eta\in(-\infty,0)$). An \emph{admissible
diamond} is a triple $D=(\eta_i,\eta_f,R_{\mathrm{conf}})$ such that
\begin{enumerate}
\item $0<\eta_i<\eta_f<\infty$ (so that $a$ is smooth and strictly
positive on $\overline{D}$),
\item $R_{\mathrm{conf}}>0$,
\item $R_{\mathrm{conf}}\le(\eta_f-\eta_i)/2$ (the standard
``radius~$=$~half conformal-time interval'' diamond convention).
\end{enumerate}
\end{definition}

\begin{definition}[Source category]\label{def:Diam}
Let $\Diam$ be the category whose objects are admissible diamonds
(Definition~\ref{def:admissible}) and whose morphisms are inclusions
$D'\hookrightarrow D$, i.e.\
$D'\le D$ iff
$\eta_i\le\eta_i'$, $\eta_f'\le\eta_f$, $R_{\mathrm{conf}}'\le R_{\mathrm{conf}}$.
\end{definition}

\begin{remark}\label{rem:poset}
$\Diam$ is a thin category (poset). Reflexivity, antisymmetry up to
equality of triples, and transitivity are verified componentwise on
the open cone of admissible parameters; this is mechanically checked
in \texttt{DME\_functor.py} Block~1 on $\ge 2000$ random admissible
chains.
\end{remark}

\section{Three candidate target categories}\label{sec:target}

\begin{itemize}
\item[$\mathbf{Cat_1}$:] W$^*$-algebras with W$^*$-isomorphisms.
\emph{Rejected.} Inclusions $D'\subset D$ in QFT do not generally
induce isomorphisms but injections (isotony of the Haag--Kastler
net).
\item[$\mathbf{Cat_2}$:] W$^*$-dynamical systems with covariant
W$^*$-isomorphisms (Connes 1973). \emph{Rejected.} Same reason; in
addition, the modular flow $\sigma^D$ on $\NN(D)$ depends on $D$ in
such a way that no covariant isomorphism between
$(\NN(D),\sigma^D)$ and $(\NN(D'),\sigma^{D'})$ is induced by the
inclusion (this is precisely the absence of a state-preserving
conditional expectation in general; see Takesaki~\cite{Takesaki1972}).
\item[$\mathbf{Cat_3}\equiv\WCatinj$:] W$^*$-algebras with
\emph{unital injective normal $*$-homomorphisms} as morphisms; the
Brunetti--Fredenhagen--Verch (CMP~237 (2003) 31) target category for
locally covariant QFT. \emph{Selected.} It is the categorical
formalisation of isotony of nets of local algebras and supports the
notion of a covariant quantum field as a \emph{natural transformation}
between functors out of $\Diam$.
\end{itemize}

\section{The DME assignment as a partial functor}\label{sec:functor}

\begin{theorem}[Functoriality of $\NN$, conditional]\label{thm:N-functor}
Assume the modular-compatibility hypothesis $(\star)$:\ for each
inclusion $D'\subset D$ in $\Diam$, there exists a faithful normal
state-preserving conditional expectation
$E_{D',D}:\NN(D)\to\NN(D')$ (equivalently, by Takesaki
\cite{Takesaki1972}, $\NN(D')\subset\NN(D)$ is a Takesaki sub-pair
for $\omega_D$). Then
\[
\NN \;:\; \Diam \longrightarrow \WCatinj,
\qquad
D \mapsto \NN(D), \qquad
(D'\hookrightarrow D)\mapsto\bigl(\iota_{D',D}:\NN(D')\hookrightarrow\NN(D)\bigr),
\]
is a covariant functor.
\end{theorem}

\begin{proof}
On objects: by Theorem~3.5/3.6 and Corollary~3.7 of
\texttt{frw\_note.tex}~\cite{Remondiere2026}, each admissible $D$
yields a type II$_\infty$ factor $\NN(D)$. On morphisms: the
inclusion of test functions $C_c^\infty(D')\subset C_c^\infty(D)$
induces $\AA(D')_\FRW\hookrightarrow\AA(D)_\FRW$ as W$^*$-subalgebras
(isotony of the Haag--Kastler net for the conformally coupled
massless scalar). Crossing with the modular flow $\sigma^D$
restricted to $\AA(D')_\FRW$, hypothesis $(\star)$ ensures that
$\sigma^D|_{\AA(D')_\FRW}=\sigma^{D'}$, so the crossed product
$\AA(D')_\FRW\rtimes_{\sigma^{D'}}\R$ injects into
$\AA(D)_\FRW\rtimes_{\sigma^D}\R$ as a normal unital
$*$-homomorphism. Functoriality (composition and identity) is
inherited from the inclusions of the underlying test-function
spaces.
\end{proof}

\begin{remark}[Status of $(\star)$]\label{rem:star}
$(\star)$ holds in the conformal-pullback unitary equivalence
$\AA(D)_\FRW\simeq\AA(M_D)_\Mink$ on Minkowski diamonds of the
Hislop--Longo type, since the Minkowski net satisfies the
Bisognano--Wichmann property and inclusions of double cones admit
state-preserving conditional expectations~\cite{Longo1979,
Buchholz1974}. For the FRW-side this transfers via the unitary
$U$ of \texttt{frw\_note} Prop.~3.4, in the conformally coupled
massless case treated by~\cite{Remondiere2026}. Beyond this case
$(\star)$ is open.
\end{remark}

\begin{proposition}[Non-functoriality of $K$]\label{prop:K-not-functor}
The assignment $D\mapsto K(D):=-\log\Delta_{\omega_D}$ does
\emph{not} extend to a functor $\Diam\to(\text{self-adjoint
operators on a fixed Hilbert space, with intertwiners})$, because
the modular operator of an inclusion $\NN(D')\subset\NN(D)$ is not
in general the restriction of $\Delta_{\omega_D}$ (Takesaki's
counterexample~\cite{Takesaki1972}). What \emph{is} functorial is
the modular automorphism \emph{group}: by Connes's cocycle theorem
(\cite{Connes1973}, Th\'eor\`eme~1.2.4), there exists a
$\sigma^{D'}$-cocycle $u_{D',D}(t)\in\NN(D')$ with
\[
\sigma^D_t \circ \iota_{D',D} \;=\; \mathrm{Ad}\bigl(u_{D',D}(t)\bigr)\circ \iota_{D',D}\circ\sigma^{D'}_t,
\]
and the assignment $D\mapsto[\sigma^D]$ (the outer-modular class) is
functorial up to inner automorphisms.
\end{proposition}

\begin{proof}
Sympy verification of the linearised cocycle identity
$u_{D',D}(s+t)=u_{D',D}(s)\,\sigma^{D'}_s(u_{D',D}(t))$ on three
independent profiles is recorded in \texttt{DME\_functor.py}
Block~3, and matches the $O(m^2)$ check of \texttt{frw\_note.tex}
Proposition~4.2.
\end{proof}

\section{The KDC as a candidate natural transformation}\label{sec:nat}

Let $\R^*_+$ denote the discrete (only-identities) category of
positive reals. Define two functors $\Diam\to\R^*_+$:
\[
F_C(D) \;:=\; \lim_{t\to\infty}\frac{1}{C_k(D;t)}\frac{\mathrm{d}C_k(D;t)}{\mathrm{d}t},
\qquad
F_R^{-1}(D) \;:=\; \frac{1}{\Rprop(\eta_c(D))}.
\]

\begin{theorem}[KDC as a non-natural transformation]\label{thm:KDC-nat}
Pointwise on each object $D\in\Diam$, the Krylov--Diameter
Correspondence (\texttt{krylov\_diameter} Thm.~4) gives
$F_C(D)=F_R^{-1}(D)$ in the late-modular-time regime under the
universal-slope hypothesis $\lambda_L^{\mathrm{mod}}=2\pi$ of
CMPT~\cite{CMPT24}. Sympy verification of the chain-rule identity
$(1/C_k)\,\mathrm{d}C_k/\mathrm{d}\tau=1/\Rprop$ on the
Parker--Cao saturating profile $C_k(s)=\sinh(\pi s)^2/\pi^2$ is
recorded in \texttt{DME\_functor.py} Block~4.

However, the assignment $\phi_D := \mathrm{id}_{F_C(D)}$ does
\emph{not} define a natural transformation $\phi:F_C\Rightarrow
F_R^{-1}$, because the naturality square
\[
\begin{tikzcd}
F_C(D') \arrow[r, "\phi_{D'}"] \arrow[d, "F_C(\iota_{D',D})"'] & F_R^{-1}(D') \arrow[d, "F_R^{-1}(\iota_{D',D})"] \\
F_C(D) \arrow[r, "\phi_D"'] & F_R^{-1}(D)
\end{tikzcd}
\]
need not commute: the right vertical is multiplication by
$\Rprop(D')/\Rprop(D)$ (a contraction), and is well-defined; but
the left vertical requires the Lanczos coefficients of the spread
complexity to be preserved by the Connes--Takesaki conditional
expectation $E_{D',D}:\NN(D)\to\NN(D')$, which is open in general
(\texttt{krylov\_diameter} \S6, Open Question regarding
projectability of Lanczos sequences).
\end{theorem}

\begin{proof}[Proof sketch]
Pointwise identity: $C_k(D;s)\sim e^{2\pi s}/(4\pi^2)$ at late
modular time on the Parker--Cao saturating profile, hence
$(1/C_k)\,\mathrm{d}C_k/\mathrm{d}s\to 2\pi$. The
Casini--Huerta--Myers relation $\mathrm{d}s/\mathrm{d}\tau=1/(2\pi
\Rprop(\eta_c))$ on the central worldline, plus the chain rule,
gives $(1/C_k)\,\mathrm{d}C_k/\mathrm{d}\tau=1/\Rprop(\eta_c)$.
This is checked in \texttt{DME\_functor.py} Block~4.

Failure of naturality: under $\iota_{D',D}$, the seed
$|\mathrm{seed}\rangle\in L^2(\NN(D'),\mathrm{tr}_{\NN(D')})$
does not in general image to a Krylov-cyclic vector in
$L^2(\NN(D),\mathrm{tr}_{\NN(D)})$; even if it does, the Lanczos
coefficients $b_n^{D'}$ on $\NN(D')$ need not equal the
restriction $b_n^D|_{\mathrm{Krylov}(D')}$ of the Lanczos
coefficients on $\NN(D)$. The conditional expectation
$E_{D',D}$ projects vectors but does not in general intertwine
the Lanczos recursion (this requires $K_{\NN(D)}$ to commute with
$E_{D',D}$, which is the modular-invariance hypothesis of
Connes--Takesaki~\cite{ConnesTakesaki1977} and is not
established for the FRW crossed product beyond the
conformal-pullback regime).
\end{proof}

\section{Limits and colimits}\label{sec:limits}

\begin{proposition}[Colimit: growing diamonds]\label{prop:colim}
The diagram $D\subset D'\subset\cdots$ of nested admissible
diamonds, ordered by inclusion, has a colimit in $\WCatinj$
given by the inductive limit
\(
\NN_{\mathrm{global}} := \varinjlim_{D\in\Diam}\NN(D),
\)
which is the global FRW crossed-product algebra. This is well-defined
under hypothesis $(\star)$ of Theorem~\ref{thm:N-functor} and is
type II$_\infty$ if all $\NN(D)$ are.
\end{proposition}

\begin{question}[Limit: Big-Bang shrinking]\label{q:limit-BB}
The diagram $D\supset D'\supset\cdots$ of shrinking admissible
diamonds with $\eta_i\to 0^+$ is the categorical realisation of
\texttt{frw\_note} Open Question~4.4. Its limit in $\WCatinj$ is
formally
\(
\NN_{\mathrm{BB}} := \varprojlim_{D\,:\,\eta_i\to 0^+}\NN(D),
\)
but the IR singularity of the conformal vacuum at $\eta=0$
(\texttt{frw\_note} Q.4.4 (b)--(c)) means the projective limit may
fail to admit a cyclic separating vector. The type-classification of
$\NN_{\mathrm{BB}}$, if it exists, is open.
\end{question}

\section{Verdict}

\begin{remark}[Summary]\label{rem:verdict}
The Diamond Master Equation is a \emph{partial} functor:
\begin{itemize}
\item $\NN:\Diam\to\WCatinj$ is a functor under hypothesis $(\star)$,
which is provable in the conformal-pullback regime of
\texttt{frw\_note} but open in general.
\item $K:D\mapsto-\log\Delta_{\omega_D}$ is \emph{not} functorial as
an operator-valued field; only the outer-modular class $[\sigma^D]$
is functorial up to inner perturbations (Connes 1973).
\item The Krylov--Diameter Correspondence is a \emph{pointwise}
identity on each $D\in\Diam$ but \emph{not} a natural transformation
between functors $\Diam\to\R^*_+$, because Lanczos coefficients are
not preserved by the conditional expectation
$E_{D',D}:\NN(D)\to\NN(D')$.
\item The colimit (growing $D$) is the global FRW algebra; the limit
(shrinking to BB) is the open question of \texttt{frw\_note}~4.4.
\end{itemize}
No new theorem is established. The construction is an honest
categorical re-packaging of three pre-existing theorems
(\texttt{frw\_note} 3.5/3.6, \texttt{krylov\_diameter} Thm.~4,
\texttt{frw\_note} Prop.~4.2). The natural target for publication of
this re-packaging, in the absence of new content, would be a
specialist functorial-QFT venue (e.g.\ JPAA or Theory and
Applications of Categories) only \emph{after} hypothesis $(\star)$
and the Lanczos-projection question are resolved into actual
theorems. As of 2026-05-02 this has not been done, and the present
note is a research-direction marker, not a publishable result.
\end{remark}

\bigskip
\noindent\emph{Status (2026-05-02):} speculative categorical
synthesis; sympy-internal-consistency only; NO new theorem; NO arXiv
claim; NO push.

\begin{thebibliography}{99}
\bibitem{Connes1973} A.~Connes, \emph{Une classification des
facteurs de type III}, Ann.\ Sci.\ \'Ec.\ Norm.\ Sup\'er.\ (4)
\textbf{6} (1973), 133--252.
\bibitem{ConnesTakesaki1977} A.~Connes and M.~Takesaki, \emph{The
flow of weights on factors of type III}, T\^ohoku Math.\ J.\
\textbf{29} (1977), 473--575.
\bibitem{Takesaki1972} M.~Takesaki, \emph{Conditional expectations
in von Neumann algebras}, J.\ Funct.\ Anal.\ \textbf{9} (1972),
306--321.
\bibitem{Buchholz1974} D.~Buchholz, \emph{Product states for local
algebras}, Comm.\ Math.\ Phys.\ \textbf{36} (1974), 287--304.
\bibitem{Longo1979} R.~Longo, \emph{Notes on algebraic invariants
for non-commutative dynamical systems}, Comm.\ Math.\ Phys.\
\textbf{69} (1979), 195--207.
\bibitem{DHR1971} S.~Doplicher, R.~Haag, J.E.~Roberts,
\emph{Local observables and particle statistics I},
Comm.\ Math.\ Phys.\ \textbf{23} (1971), 199--230;
\emph{II}, Comm.\ Math.\ Phys.\ \textbf{35} (1974), 49--85.
\bibitem{BFV2003} R.~Brunetti, K.~Fredenhagen, R.~Verch, \emph{The
generally covariant locality principle -- a new paradigm for local
quantum physics}, Comm.\ Math.\ Phys.\ \textbf{237} (2003), 31--68.
\bibitem{BuchholzWichmann1986} D.~Buchholz, E.H.~Wichmann,
\emph{Causal independence and the energy-level density of states in
local quantum field theory}, Comm.\ Math.\ Phys.\ \textbf{106}
(1986), 321--344.
\bibitem{Longo1989} R.~Longo, \emph{Index of subfactors and
statistics of quantum fields, I}, Comm.\ Math.\ Phys.\ \textbf{126}
(1989), 217--247; \emph{II}, Comm.\ Math.\ Phys.\ \textbf{130}
(1990), 285--309.
\bibitem{CMPT24} P.~Caputa, J.~M.~Magan, D.~Patramanis, E.~Tonni,
\emph{Krylov complexity for von Neumann algebras with type II$_1$ /
III$_1$ factors}, 2024 (referenced in
\texttt{krylov\_diameter.tex}).
\bibitem{Remondiere2026} K.~Remondi\`ere, \emph{Type II$_\infty$ for
the gravitational algebra of a comoving observer in
radiation-dominated FRW via conformal pullback}, manuscript
(\texttt{frw\_note.tex}, commit \texttt{aa9a3ee} of
\texttt{crossed-cosmos}, 2026-05-02).
\end{thebibliography}

\end{document}
